% Generated draft: requires mathematical and source review.
\documentclass[11pt]{article}
\usepackage[margin=25mm]{geometry}
\usepackage{fontspec}
\setmainfont{Latin Modern Roman}
\usepackage{amsmath,amssymb,mathtools,unicode-math}
\setmathfont{Latin Modern Math}
\usepackage{enumitem,needspace,xurl,hyperref}
\setlength{\parindent}{0pt}
\setlength{\parskip}{5pt}
\setlength{\emergencystretch}{3em}
\newcommand{\N}{\mathbb{N}}
\newcommand{\Z}{\mathbb{Z}}
\newcommand{\Q}{\mathbb{Q}}
\newcommand{\R}{\mathbb{R}}
\newcommand{\C}{\mathbb{C}}
\newcommand{\F}{\mathbb{F}}
\newcommand{\A}{\mathbb{A}}
\newcommand{\PP}{\mathbb{P}}
\newcommand{\E}{\mathbb{E}}
\newcommand{\eps}{\varepsilon}
\newcommand{\dd}{\,\mathrm{d}}
\begin{document}
\section*{UnsolvedMath ID: 20001851}
\textbf{Source problem code:} catalogue entry
\subsection{Definitions and mathematical statement}
The source, attributed to Streinu, asks whether the space of non-intersecting configurations of an equilateral open arm in \(\mathbb R^3\) is connected. Its damaged length condition is explicitly reconstructed by the supplied verification as \(\ell_1=\cdots=\ell_n\). The reference ``same as the previous question'' has no preceding question in this record.

In the strict zero-thickness interpretation, a configuration is a tuple \((p_0,\ldots,p_n)\) with \(\|p_i-p_{i-1}\|=\ell_i>0\). Nonadjacent closed bars are disjoint, and adjacent bars intersect only at their common endpoint. An open arm has no endpoint identification. For the equilateral question rescale to \(\ell_i=1\), fix \(p_0=0\), and retain the topology inherited from vertex coordinates. The question is whether every two such configurations can be joined by a continuous path of strict non-intersecting configurations preserving all bar lengths. In this finite semialgebraic model connectedness and path connectedness agree. The anchoring and collision convention are explicit reconstructions.

The aid reports that the arbitrary-length equilateral problem appears to remain open, with published work implying connectivity for at most five unit bars. It gives a sufficient certificate for a particular chain, allowing arbitrary positive lengths. For a unit axis \(a\), let
\[
 I_i(a)=\left[\min\{a\cdot p_{i-1},a\cdot p_i\},
                  \max\{a\cdot p_{i-1},a\cdot p_i\}\right].
\]
The certificate is
\[
 I_i(a)\cap I_j(a)=\varnothing\qquad(|i-j|\geq2).
\]
These are closed intervals, so the separation is strict. The reported lemma asserts that such a chain has a simple orthogonal projection onto a suitable nearby generic plane and therefore can be straightened, using the theorem attributed to Biedl et al. The axis/plane choice may be perturbed to avoid projection degeneracies while retaining the finite strict separations. The certificate persists under sufficiently small perturbations of the vertices and axis.

The aid further reports that the certificate is automatic for every strict simple three-bar chain of arbitrary positive lengths: its sole nonadjacent pair consists of disjoint compact segments and admits strict scalar separation. This is a sufficient unlockability statement for those chains. It does not show that every equilateral chain of arbitrary length has this certificate, and failure to find a certifying axis is not proof of locking.
\subsection{Short English statement}
Streinu's question asks whether all non-intersecting equilateral open chains in three-space lie in one configuration-space component. The supplied report leaves arbitrary length open and states small-chain results. Its additional finite certificate uses one axis separating every pair of nonadjacent bars in projection, which guarantees a simple planar projection and hence straightening.
\subsection{Sources}
\begin{itemize}
\item Local source record: catalogue entry, catalogue entry; canonical \texttt{aim-geometry-notes.json}, \texttt{notes[188]}, source item 2. Title: ``A one-coordinate unlockability certificate for open chains in three-space.''
\item AIM workshop ``Configuration spaces of linkages,'' problem 2, attributed to Streinu; supplied PDF source reference \url{https://aimath.org/pastworkshops/linkagesproblems.pdf}. This draft uses the supplied text and does not access or compile a PDF. Original remarks and literature fields are empty.
\item Supplied verification says the canonical JSON lost the letter ell in the equal-length condition. The research aid attributes the simple-projection straightening theorem to Biedl et al. and invokes published connectivity for at most five unit bars, without full bibliographic entries.
\item Registry dated 14 August 2026: ``partially\_solved,'' ``PARTIAL-PROGRESS,'' statement status ``exact.'' Attempt 1 is partial, research-run status open, rigor label ``proved,'' novelty confidence low. The machine-generated, non-peer-reviewed aid expressly requires independent verification before citation.
\end{itemize}
\end{document}
